\begin{thebibliography}{99}
	\RaggedRight
	
	% 参考文献间距压缩
	\setlength{\itemsep}{0pt}
	\setlength{\parskip}{0pt}
	\setlength{\parsep}{0pt}
	
	% 单独设置参考文献行距
	\linespread{1.5}\selectfont
	
	\setlength{\emergencystretch}{2em}
	
	\bibitem{ref1}
	王继业，周碧玉，张法，等，数据中心能耗模型及能效算法综述，
	计算机研究与发展，56(8)：1587--1603，2019。	
	\bibitem{ref2}
	丁肇豪，曹雨洁，张素芳，等，能源互联网背景下数据中心与电力系统协同优化（一）：数据中心能耗模型，
	中国电机工程学报，42(9)：3161--3176，2022。	
	\bibitem{ref3}
	吕佳炜，张沈习，程浩忠，等，集成数据中心的综合能源系统能量流--数据流协同规划综述及展望，
	中国电机工程学报，41(16)：5500--5521，2021。
	
	\bibitem{ref4}
	Wickramasuriya S L, Athanasopoulos G, Hyndman R J,
	Optimal Forecast Reconciliation for Hierarchical and Grouped Time Series Through Trace Minimization,
	Journal of the American Statistical Association, 114(526): 804--819, 2019.
	
	\bibitem{ref5}
	Radovanovic A, Koningstein R, Schneider I, et al.,
	Carbon-Aware Computing for Datacenters,
	IEEE Transactions on Power Systems, 38(2): 1270--1280, 2023.
	
	\bibitem{ref6}
	Lindberg J, Lesieutre B C, Roald L A,
	Using Geographic Load Shifting to Reduce Carbon Emissions,
	Electric Power Systems Research, 212: 108586, 2022.
	
	\bibitem{ref7}
	Guo X, Che Y, Zheng Z, et al.,
	Multi-timescale Optimization Scheduling of Interconnected Data Centers Based on Model Predictive Control,
	Frontiers in Energy, 18(1): 28--41, 2024.
	
	\bibitem{ref8}
	Liu W, Yan Y, Sun Y, et al.,
	Online Job Scheduling Scheme for Low-carbon Data Center Operation: An Information and Energy Nexus Perspective,
	Applied Energy, 338: 120918, 2023.
	
	\bibitem{ref9}
	向梓旸，黄淳驿，李康平，等，考虑负载时空转移特性的数据中心集群算电协同优化调度，
	电力信息与通信技术，24(3)：33--46，2026。
	
	\bibitem{ref10}
	Wang K, Ye L, Yang S, et al.,
	A Hierarchical Dispatch Strategy of Hybrid Energy Storage System in Internet Data Center with Model Predictive Control,
	Applied Energy, 331: 120414, 2023.

	\bibitem{ref11}
	DeepSeek, DeepSeek-v4-flash-0731, 深度求索 (DeepSeek), 使用日期：2026-08-07--2026-08-10.
	
\end{thebibliography}
